\documentclass[11pt]{article}
\usepackage{ochamath}
\subjectcolor{2F5DA8}
\setsheet{線形代数}{ベクトル空間と部分空間　演習}{https://math.ochanote.com/linear-algebra/vector-spaces/}
\namelinetrue
\begin{document}
\makesheethead{問題}

\problem[部分空間の判定]
$\mathbb R^2$ の集合 $W_1=\{(x,y):y=3x\}$、$W_2=\{(x,y):y=3x+1\}$ は部分空間か。理由を述べよ。
\workspace{38mm}

\problem[生成の範囲]
$\boldsymbol u=(1,0,1)^{\mathsf T}$、$\boldsymbol v=(0,1,1)^{\mathsf T}$ から $(2,-1,1)^{\mathsf T}$ と $(0,0,1)^{\mathsf T}$ を作れるか。
\workspace{38mm}
\newpage

\problem[独立性]
$\boldsymbol a=(1,2)^{\mathsf T}$、$\boldsymbol b=(2,4)^{\mathsf T}$ の独立性を判定せよ。次に $(1,0)^{\mathsf T},(0,1)^{\mathsf T}$ の組と比べよ。
\workspace{38mm}

\problem[出力制約]
出力 $\boldsymbol y=s(1,0,1)^{\mathsf T}+t(0,1,1)^{\mathsf T}$ の成分の関係を求めよ。$(3,2,5)^{\mathsf T}$ は到達可能か。
\workspace{38mm}
\end{document}
